% Aufgaben für den LaTeX-Kurs (2017-10-06)
% Autor: Simon May
% Datum: 2017-10-04

\documentclass[a4paper, 12pt, headsepline, parskip=half-, headlines=3.1]{scrartcl}

\usepackage{amsmath}
\usepackage{mathtools}

\usepackage{unicode-math}
\usepackage{polyglossia}
\setdefaultlanguage{german}

\usepackage{selnolig}
\usepackage{microtype}
\usepackage{scrlayer-scrpage}
\usepackage{lastpage}
\usepackage[useregional]{datetime2}

\usepackage{csquotes}
\usepackage{siunitx}
\usepackage{enumitem}
\usepackage{xcolor}
\usepackage{wrapfig}
\usepackage{graphicx}
\usepackage{grffile}
\usepackage{float}
\usepackage{listingsutf8}
\usepackage[backend=biber, style=verbose]{biblatex}
\usepackage{xfrac}
\usepackage{subcaption}

\usepackage[unicode, pdfusetitle]{hyperref}
\usepackage{cleveref}

% --- Paket-Einstellungen
% -- latex
\widowpenalty=9000
\clubpenalty=9000
\removenolimits{\int}
% -- koma
\setkomafont{pageheadfoot}{\normalfont}
\addtokomafont{pagefoot}{\small}
% -- scrlayer-scrpage
\ihead[]{\textbf{Einführung in \LaTeX\\Aufgaben}}
\chead[]{\includegraphics[height=2.8\baselineskip]{res/logo.pdf}}
\ohead[]{Simon May\\Fachschaft Physik}
\cfoot[Seite: {\pagemark} von \pageref{LastPage}]{Seite: \pagemark\ von \pageref{LastPage}}
\pagestyle{scrheadings}
% -- biblatex
\addbibresource{res/literatur.bib}
% -- csquotes
\MakeOuterQuote{"}
% -- siunitx
\sisetup{
	locale=DE,
	output-product=\cdot,
	separate-uncertainty,
	quotient-mode=fraction,
	per-mode=fraction,
	fraction-function=\sfrac,
}
% -- listings
\colorlet{keywordcolor}{purple!60!black}
\colorlet{commentcolor}{green!40!black}
\colorlet{texcsicolor}{blue}
\colorlet{texcsiicolor}{blue!50!black}
\colorlet{bracescolor}{orange!70!black}
\colorlet{dollarcolor}{red}
\lstdefinestyle{mystandard}{
	% Stelle UTF8-Input mit Latin1 dar
	%inputencoding=utf8/latin1,
	% Zeilennummern links
	%numbers=left,
	% Darstellung der Zeilennummern
	numberstyle=\ttfamily,
	% Rahmen (links)
	%frame=L,
	columns=flexible,
	%xleftmargin=0.5cm,
	% Default Schriftart/Größe
	basicstyle=\ttfamily,
	% Darstellung für Keywords
	keywordstyle=\color{keywordcolor},
	% Darstellung für Kommentare
	commentstyle=\color{commentcolor},
	% Darstellung für Strings
	stringstyle=\color[RGB]{ 47,   0, 255},
	% Darstellung für TeX-Befehle; Stern: mit backslash
	texcsstyle=*\color{texcsicolor},
	texcsstyle=*[2]\color{texcsiicolor},
	% Tab-Breite
	tabsize=4,
	% Leerzeichen in Strings hervorheben
	showstringspaces=false,
	% Tabs hervorheben
	showtabs=false,
	% Leerzeichen hervorheben
	showspaces=false,
	% Alle Leerzeichen behalten
	keepspaces=true,
	% "Normale" Anführungszeichen: ""
	upquote=true,
	% Lange Zeilen aufbrechen
	breaklines=true,
	% Zeilen an Whitespace aufbrechen
	breakatwhitespace=true,
}

\lstset{
	language=[LaTeX]{TeX},
	style=mystandard,
	escapechar=|,
	deletetexcs={begin, end, documentclass, usepackage},
	moretexcs=[2]{begin, end, documentclass, usepackage},
	moretexcs={
		f, latex, Alpha,
		part, chapter, subsection, subsubsection, paragraph, subparagraph, color, text, mathbb, mathcal, tableofcontents, includegraphics, cref, crefrange, autocite, footcite, printbibliography,
		SI, num, si, SIrange, SIlist, m, per, s, squared, cm, nm, MHz, mega, Hz, cubed, tothe, mA, kV
	},
	literate=
		{\$}{\textcolor{dollarcolor}{\$}}1
		{\{}{\textcolor{bracescolor}{\{}}1
		{\}}{\textcolor{bracescolor}{\}}}1,
}


% --- Befehle
\newcommand{\email}[1]{\href{mailto:#1}{\texttt{#1}}}

\newcounter{aufgabe}
\newenvironment{aufgabe}[1][]{\stepcounter{aufgabe}
\noindent\textbf{Aufgabe \arabic{aufgabe}.} #1\par
\begin{enumerate}[label=\alph*)]}{\end{enumerate}\bigskip}
\newenvironment{aufgabe*}[1][]{\stepcounter{aufgabe}
\noindent\textbf{Aufgabe \arabic{aufgabe}.} #1}{\bigskip}

\title{Einführung in \LaTeX: Aufgaben}
\author{Simon May\\\email{simon.may@uni-muenster.de}}
\date{\DTMdate{2017-10-06}}

\begin{document}

\maketitle

\begin{aufgabe}[Textformatierung]
	\item Stellt das folgende Beispiel aus der Vorlesung nach:
	\begin{itemize}
		\item \textbf{fett}
		\item \textit{kursiv}
		\item \texttt{Schreibmaschine}
		\item \underline{unterstrichen}
		\item \textsc{Kapitälchen}
	\end{itemize}
	\item Stellt diesen Satz nach:
	
	{\huge"\tiny Das ist \small ein Satz, \large der \Large diese \huge Schriftgrößen demonstriert."}
\end{aufgabe}

\begin{aufgabe*}
	Aufzählungen – Stellt das folgende Beispiel nach:
	\begin{enumerate}
		\item Physik~I
		\begin{itemize}
			\item Methodik der Physik
			\item Klassische (newtonsche) Mechanik
			\item Bewegungsgleichungen
			\item Einblick in die spezielle Relativitätstheorie
		\end{itemize}
		\item Physik~II
		\begin{enumerate}
			\item Kinetische Gastheorie
			\begin{enumerate}
				\item Maxwell-Boltzmann-Verteilung
				\item Ideales Gas
				\item Modelle realer Gase
			\end{enumerate}
			\item Klassische Thermodynamik
			\item Elektrostatik
			\item Magnetostatik
		\end{enumerate}
		\item Physik~III
		\begin{description}
			\item[Klassische Elektrodynamik] Maxwell-Gleichungen, Wechselstrom, Licht
			\begin{enumerate}
				\item Elektrodynamik im Vakuum
				\item Elektrodynamik in Materie
			\end{enumerate}
			\item[Ausbreitung von Wellen] Wellen auf Kabeln, Lichtausbreitung im Vakuum
			\item[Quantenphysik] Einblick in die Quantenmechanik
			\item[Spezielle Relativitätstheorie] Schnell bewegte Bezugssysteme
		\end{description}
	\end{enumerate}
\end{aufgabe*}

\begin{aufgabe}[Formeln]
	\item Setzt die folgenden Formeln mit \LaTeX:
	\begin{equation}
		E = mc^2
	\end{equation}	
	\begin{equation}
		\label{eq:newton}
		F = -G ⋅ \frac{mM}{r^2}
	\end{equation}	
	\begin{equation}
		C = \frac{Q}{U}
	\end{equation}	
	\begin{equation}
		U = U_0 ⋅ e^{-\frac{t}{RC}}
	\end{equation}	
	\begin{equation}
		\label{eq:circle_circumference}
		U = 2πr
	\end{equation}
	\begin{equation}\label{eq:circle_area}
		A = πr^2 = π\frac{d^2}{4}
	\end{equation}
	\begin{equation}
		r = \sqrt{\frac{(s-a)(s-b)(s-c)}{s}}
	\end{equation}
	\item Fügt eine Formel in den Fließtext ein:
	
	"Der Satz von Pythagoras, $a^2 + b^2 = c^2$, beschreibt das Verhältnis von Längen in einem rechtwinkligen Dreieck und ist eine sehr bekannte mathematische Aussage."
	\item Setzt folgende Herleitung mit \LaTeX:
	\begin{align*}
		(a + b)^3
		&= (a + b)^2 (a + b) \\
		&= (a^2 + 2ab + b^2) (a + b) \\
		&= \left[(a^3 + 2 a^2 b + a b^2) + (a^2 b + 2 a b^2 + b^3)\right] \\
		&= a^3 + 3 a^2 b + 3 a b^2 + b^3
	\end{align*}
	\item Versucht euch als Nächstes an diesen Formeln:
	\begin{align*}
		0 &= \frac{1}{π} ∫_0^{2π} \cos(nx) \sin(mx) \, \mathrm{d}x\\
		\dot{Q} &= - λA \frac{ΔT}{L}\\
		I &= ∫_A \vec{j} ⋅ \mathrm{d}\vec{A}\\
		\vec{F} &= - \frac{q_1 q_2}{4πε_0} \frac{\vec{r}_1 - \vec{r}_2}{|\vec{r}_1 - \vec{r}_2|^3} \\
		A &= \frac{\tan(x) \sin(x)}{π^2 ρ_{\text{neu}}} ⋅ ρ_{\text{alt}} ⋅ ∫_1^X \exp\!\left(\frac{m v^2}{2\ln(x)}\right) \, \mathrm{d}x
	\end{align*}
	\item Setzt die Maxwell-Gleichungen in Materie mit \LaTeX:
	\begin{align}
		\label{eq:max1}
		∇ ⋅ \vec{D} &= ρ_\text{f} \\
		∇ ⋅ \vec{B} &= 0 \\
		∇ × \vec{E} &= -\frac{∂\vec{B}}{∂t} \\
		\label{eq:max4}
		∇ × \vec{H} &= \vec{j}_\text{f} + \frac{∂\vec{D}}{∂t}
	\end{align}
	
	\emph{Erinnerung:} \lstinline!\nabla! ergibt $∇$, \lstinline!\partial! ergibt $∂$.
\end{aufgabe}

\clearpage

\begin{aufgabe}[Verweise]
	\item Versucht, einen automatischen Verweis auf \cref{eq:newton} zu erzeugen.
	\item Erzeugt den folgenden Text mit automatischen Verweisen:\\
	"\cref{eq:newton,eq:circle_circumference,eq:circle_area}".
	\item Erzeugt den folgenden Text mit automatischen Verweisen:\\
	"\crefrange{eq:max1}{eq:max4}".
	
	\emph{Hinweis:} Einen Bereich von Verweisen kann man mit\\
	\lstinline!\crefrange{<label1>}{<label2>}! erzeugen.
\end{aufgabe}

\begin{aufgabe}[Tabellen – Reproduziert die folgenden Tabellen.]
	\item 
	\begin{tabular}{|c c || c|}
		\hline
		1 & 2 & 3 \\ \hline\hline
		4 & 5 & 6 \\
		7 & 8 & 9 \\ \hline
	\end{tabular}
	\item
	\begin{tabular}{l | S S S S | S}
		& \multicolumn{4}{c|}{\textbf{Augenfarbe}} & \\
		\textbf{Haarfarbe}	& \textbf{Braun}	& \textbf{Blau} & \textbf{Grau}
			& \textbf{Grün}		& \textbf{Summe} \\ \hline
		\textbf{Braun}		& 119	& 84	& 54	& 29	& 286	\\
		\textbf{Blond}		& 7		& 94	& 10	& 16	& 127	\\
		\textbf{Schwarz}	& 68	& 20	& 15	& 5		& 108	\\
		\textbf{Rot}		& 26	& 17	& 14	& 14	& 71	\\ \hline
		\textbf{Summe}		& 220	& 215	& 93	& 64	& 592
	\end{tabular}
	
	\emph{Hinweis:} Mit dem Befehl\\
	\lstinline!\multicolumn{<spaltenzahl>}{<neue_spalte>}{Text...}! lassen sich Spalten zusammenfassen.

	\emph{Hinweis:} Zur einfacheren Eingabe gibt es z.\,B.\ in TeXStudio die Option "Assistenten $→$ Tabellen-Assistent…"
	\item Cyclohexan:
	
	\begin{tabular}{S S}
		\hline
		{Masse / \si{\g}}	& {Volumenänderung / \si{\ml}} \\
		\hline
		0,15		& 55+-2 \\
		0,18		& 59+-1 \\
		0,19		& 58+-1 \\
		0,19		& 55+-1 \\
		0,2			& 61+-3 \\
		\hline
	\end{tabular}
		
	\emph{Hinweis:} Verwendet \lstinline!S!-Spalten!
	
	\clearpage
	\item ~
	
	\begin{table}[H]
		\centering
		\caption{Gemessene Massen und daraus resultierende Volumenänderung bei der
			Bestimmung der molaren Masse von Ethanol mit der Dampfdichtemethode.}
		\begin{tabular}{S S}
			\hline
			{Masse / \si{\g}}	& {Volumenänderung / \si{\ml}} \\
			\hline
			0,06		& 55+-3 \\
			0,1			& 60+-2 \\
			0,11		& 58+-1 \\
			0,07		& 54+-1 \\
			0,08		& 52+-1 \\
			\hline
		\end{tabular}
		\label{tab:ethanol}
	\end{table}
\end{aufgabe}

\begin{aufgabe}[Abbildungen]
	\newcommand{\kursangle}{240}
	\item Fügt folgendes Bild (\texttt{answer.jpg}) in euer Dokument ein:
	
	\includegraphics[width=5cm, angle=\kursangle]{res/answer.jpg}
	
	\emph{Hinweis:} Der Winkel ist \SI{\kursangle}{\degree}.
	\medskip
	\item Erstellt die folgende Abbildung mit Unterabbildungen\\
	(zweites Bild: \texttt{8bit\_mario.png}):
	
	\begin{figure}[H]
		\centering
		\begin{subfigure}{0.45\textwidth}
		\centering
		\includegraphics[width=0.65\textwidth]{res/answer.jpg}
		\caption{Antwort auf die ultimative Frage nach dem Leben, dem Universum und dem ganzen Rest.}
		\end{subfigure}
		\hfill
		\begin{subfigure}{0.45\textwidth}
		\centering
		\includegraphics[width=0.65\textwidth]{res/8bit_mario.png}
		\caption{8bit-Mario sucht die Prinzessin.}
		\end{subfigure}
		\caption{Zwei Beispielbilder.}
	\end{figure}
	\item Erzeugt den folgenden Text mit eingebettetem Bild (\texttt{logo.pdf}):
\end{aufgabe}

\vspace{-2ex}
\enlargethispage{\baselineskip}

Die Auszeichnung der Fachschaft ging dieses Jahr an Herrn Prof.\ Dr.~Alfons Khoukaz für die Vorlesungsreihe Kern- und Teilchenphysik, welche aus einer Einführungsvorlesung im 5.~Semester und einer fortgeschrittenen Vorlesung für Masterstudenten besteht.

\begin{wrapfigure}{r}{0cm}
	\centering
	\includegraphics[width=9cm]{res/logo.pdf}
	\caption{Das Logo der Fachschaft Physik der Westfälischen Wilhelms-Universität Münster.}
\end{wrapfigure}
Da es am Fachbereich viele sehr gute Dozenten gibt, fiel die Entscheidung auch dieses Jahr wieder nicht leicht.
Das Lehrpreiskomitee musste viele Kriterien gegeneinander abwägen und hat sich dieses Jahr erstmals für eine Master-Veranstaltung entschieden.
Ebenso wurde erstmals eine ganze Vorlesungsreihe ausgezeichnet.
Der Preisträger überzeugte durch über Jahre hinweg konstant herausragende Evaluationsergebnisse sowie durch eine nahezu perfekt ausgefallene Evaluation in der diesjährigen Vorlesung "Kern- und Teilchenphysik~II".

\enlargethispage{\baselineskip}

Herr Prof.\ Dr.~Alfons Khoukaz ist Dozent am Institut für Kernphysik und Mitglied des Graduiertenkollegs "Starke und schwache Wechselwirkung".
Zu den Forschungsschwerpunkten seiner Arbeitsgruppe zählt die Mesonenproduktion in verschiedenen Stoßprozessen, was zu einem verbesserten Verständnis der starken Wechselwirkung sorgen soll.
Hierzu ist die Gruppe an einer Reihe von Experimenten beteiligt, etwa PANDA, BESIII und COSY.

\begin{aufgabe}[Erzeugt folgende Werte mit \texttt{siunitx}:]
	\item \SI{12345,67890}{\ohm}
	\item \num{1 +- 2i}
	\item \SI{.3e45}{\m\cubed}
	\item \num{1.654 x 2.34 x 8,99}
	\item \SI{1,793 +- 0,026}{\micro\F}
	\item \SI{9,81 +- 0,52}{\m\per\s\squared}
\end{aufgabe}

\begin{aufgabe*}
	Zitiert den Nolting~4\footcite{nolting4} und den Demtröder~1\footcite{demtröder1} im Dokument und erstellt das Literaturverzeichnis (\lstinline!\printbibliography!)!
\end{aufgabe*}

\begin{aufgabe*}
Baut das Dokument in der Datei \texttt{tableofcontents.pdf} nach!

\emph{Hinweis:} Das Inhaltsverzeichnis wird mit dem Befehl \lstinline!\tableofcontents! erzeugt.
\end{aufgabe*}

\begin{aufgabe*}
In der Datei \texttt{beispielprotokoll.pdf} findet ihr ein Protokoll zu einem der Versuche im Grundpraktikum. Baut es so originalgetreu wie möglich nach!

Die nötigen Bilder findet ihr im Ordner \texttt{beispielprotokoll\_bilder}.
\end{aufgabe*}

\begin{aufgabe*}
	Erzeugt die Datei \texttt{prony.pdf} mit gnuplot!

	\emph{Hinweise}:
	\begin{itemize}
		\item Die Messwerte befinden sich in der Datei \texttt{prony.csv}.
		\item Ihr dürft konstante Werte für die Fehlerbalken auf der $η$-Achse wählen!
		\item Nutzt die folgenden Formeln und Werte, um den Wirkungsgrad $η$ zu berechnen:
		\begin{align*}
			P_\text{el} &= \num{14,6} ⋅ \num{0,048} ⋅ \frac{1000}{3} \si{\W} &
			P &= 2πfFr \\
			r &= \SI{0,5}{\m} &
			η &= \frac{P}{P_\text{el}}
		\end{align*}
		\item Die Ausgleichskurve ist eine Parabel.
	\end{itemize}
\end{aufgabe*}

\end{document}
